\documentclass[12pt,a4paper]{article}
\usepackage[utf8]{inputenc}
\usepackage[spanish,es-noquoting]{babel}
\usepackage[T1]{fontenc}
\usepackage{lmodern}
\usepackage{geometry}
\geometry{top=2cm, bottom=2cm, left=2.5cm, right=2.5cm}
\usepackage{xcolor}
\usepackage{listings}
\usepackage{enumitem}
\usepackage{fancyhdr}
\usepackage{hyperref}
\usepackage{tikz}
\usetikzlibrary{arrows.meta, positioning, calc, shapes}

\definecolor{codegreen}{rgb}{0,0.6,0}
\definecolor{backcolour}{rgb}{0.95,0.95,0.92}

\lstset{
    language=C,
    backgroundcolor=\color{backcolour},
    commentstyle=\color{codegreen},
    keywordstyle=\color{blue},
    basicstyle=\ttfamily\small,
    breaklines=true,
    frame=single,
    numbers=left,
    tabsize=4,
    extendedchars=true,
    showstringspaces=false,
    literate={á}{{\'a}}1 {é}{{\'e}}1 {í}{{\'i}}1 {ó}{{\'o}}1 {ú}{{\'u}}1 {ñ}{{\~n}}1
}

\pagestyle{fancy}
\fancyhf{}
\lhead{Monitorías Sistemas Operativos 2026-I}
\rhead{Pre-Taller: Memoria Compartida}
\cfoot{\thepage}

\title{\textbf{Pre-Taller: Tablero Concurrente de Sensores}}
\author{Malak Sanchez \\ \small Monitorías de Sistemas Operativos}
\date{\today}

\begin{document}

\maketitle

\section*{Introducción}
Este ejercicio pone a prueba el manejo de memoria compartida (\texttt{SHM}) para la transferencia de información dinámica entre procesos independientes. Se busca entender cómo compartir estados globales en un espacio de memoria mapeado por el núcleo.

\section*{Problema 1: Monitoreo Concurrente}
Diseñe un programa concurrente que simule un tablero de control industrial. El programa debe recibir por argumento un número entero $\mathbf{N}$ que representa la cantidad de ciclos de actualización de cada sensor. El proceso padre (Monitor) debe crear el segmento de memoria compartida y luego disparar 3 procesos hijos (Sensores).

Cada sensor escribirá sus lecturas aleatorias en su campo asignado de la SHM exactamente $N$ veces. El Monitor imprimirá el reporte del tablero en tiempo real y solo finalizará cuando todos sus obreros hayan completado sus ciclos.

\begin{center}
\begin{tikzpicture}[node distance=2cm, font=\small, >=Stealth]
    \node[draw, rectangle, fill=blue!10, minimum width=8cm, minimum height=1cm] (shm) {MEMORIA COMPARTIDA: \texttt{[S1 | S2 | S3]}};
    
    \node[draw, circle, fill=green!10, above=of shm, xshift=-3cm] (h1) {S1};
    \node[draw, circle, fill=green!10, above=of shm] (h2) {S2};
    \node[draw, circle, fill=green!10, above=of shm, xshift=3cm] (h3) {S3};
    \node[draw, rectangle, fill=red!10, below=of shm] (padre) {MONITOR (Padre)};
    
    \draw[->, thick, dashed] (h1) -- (shm.north -| h1) node[midway, left] {Escribe ($N$)};
    \draw[->, thick, dashed] (h2) -- (shm.north) node[midway, right] {Escribe ($N$)};
    \draw[->, thick, dashed] (h3) -- (shm.north -| h3) node[midway, right] {Escribe ($N$)};
    \draw[<-, thick, color=blue!70!black] (padre) -- (shm.south) node[midway, right] {Lectura};
\end{tikzpicture}
\end{center}

\textbf{Ejemplo de ejecución:} 
\begin{lstlisting}[numbers=none]
$ ./tablero_sensores 3
[Monitor] Estado actual => S1: 45 | S2: 12 | S3: 89
[Monitor] Estado actual => S1: 67 | S2: 12 | S3: 92
[Monitor] Estado actual => S1: 67 | S2: 24 | S3: 15
[Monitor] Ciclos finalizados, memoria liberada.
\end{lstlisting}

\section*{Requerimientos}
\begin{enumerate}
    \item Implementar correctamente \texttt{shmget}, \texttt{shmat} y la creación de procesos hijos.
    \item El padre debe esperar la terminación de los procesos hijos utilizando \texttt{wait()} antes de cerrar y eliminar el segmento.
    \item Los datos deben persistir únicamente en la RAM compartida del sistema.
\end{enumerate}

\section*{Criterios de Evaluación}
\begin{itemize}
    \item \textbf{Gestión de SHM:} Creación y destrucción de recursos.
    \item \textbf{Sincronización:} Control de flujo y finalización segura de los ciclos.
    \item \textbf{Validación:} Manejo de errores en las llamadas al sistema.
\end{itemize}

\end{document}
